\section[Eleven integers with an amplified product carrier]{An eleven-integer single-product carrier at degree 648}
\label{sec:integer-product}
The new nondominated point is an ordinary polynomial in eleven unrestricted
integers, of exact degree \(648\), evaluated by \(115\) binary arithmetic
operations. It improves the recovered \(654/117\) endpoint in both degree and
cost. Three audited comparisons have degree/cost pairs \(640/118\),
\(580/119\), and \(574/122\); they are dominated by other recovered equations
and are not advertised as frontier advances. None replaces the retained
\(352/132\) low-degree or \(840/110\) low-cost endpoint.

\subsection{The complete equation}
The variables are \(j,C,z,h,r,v,w,T,c,n,s\in\mathbb Z\). All other symbols
below are polynomial abbreviations. Use the same six authenticated fixed
numerals \(\beta,K_0,M_C,M_F,A_C,B_C\), with the already shifted \(M_F\).
Their complete bytes and reconstruction recipe accompany the literal DAG.
For the primary weak-scale equation define
\begin{align*}
 J&=6j+1,&q&=\beta J+1,&W&=12w+4,\\
 X&=qW,&Y&=qv,&E&=XY,\\
 a&=E+Y,&A&=a+2,&H&=4a+3,\\
 \Delta&=a^2+H,&F&=(K_0+W)C-z(q-1),\\
 R&=(q(q-F)-C)(q^2-1)+(M_C+qM_F)J,\\
 k&=R+1+hE,&\eta&=c-Yk,\\
 D_{\rm err}&=(A_CF-q)^2+(B_CC-q)^2,\\
 B_0&=k^2(q^2-D_{\rm err})-4q^2\eta^2,&B&=B_0,\\
 \ell&=XY^2,&L&=\ell k,&D_t&=1+L(L+k),\\
 D&=3B-4,&D_g&=1+4DD_ts^2,&P&=D_tD_g,\\
 D_m&=1+\Delta c^2,&B_P&=D_mP,&U&=Arc^2,\\
 D_f&=1+\Delta U^2B_P,&Q&=\Delta^2U^2B_P,&\delta&=Q-1,\\
 V&=D_f(c-R+4UB_PT)-c,&C_y&=QV^2-1,\\
 K&=X(3c-2X)+2,&M&=2q^2,&D_c&=HM,\\
 z_c&=D_cT+MK+H(v-q^2),\\
 A_c&=(z_c^2+3D_c^2)(z_c^3-19D_c^3),\\
 S&=srHM^5,&Z&=Sn-A_c.
\end{align*}
The final polynomial is simply
\[
 \boxed{\mathcal F=\delta Z^2-S^2B_PD_fC_y.}
\]
There are no divisions, radical gates, variable exponents or unnamed
witnesses. All scalar products and powers are charged in the literal DAG.
For the primary direct equation put $\lambda=0$. The simple center replaces \(Z\) by \(S(n+\lambda V)-A_c\), where
\(\lambda=AUB_P\). The free-scale comparisons instead take
\(Y=v\), \(B=Y^2B_0-k^2q^8\), \(M=2q^3\), and replace
\(H(v-q^2)\) in \(z_c\) by \(H(v-q^3)\).

\subsection{Soundness, including the early zero-scale branch}
Unconditionally, the authentic front gives \(q\) even, \(q\equiv2\pmod3\),
\(8\mid X\), \(X\equiv2\pmod3\), and \(|X|\ge4|q|\).
Also \(H\) is nonzero odd, \(\Delta>1\), \(D_m\ge1\), and
\(D_t\ge1\). The odd integers \(D_g,P\) are nonzero. The integers
\(D_f,\delta,C_y\) cannot vanish. As
\(\pi(t)=(t^2+3)(t^3-19)\) has no rational root,
\(A_c=D_c^5\pi(z_c/D_c)\ne0\). Therefore a zero of \(\mathcal F\)
has \(S\ne0\), hence \(r,s\ne0\).

Put \(D_a=\delta C_y\). The three nonzero factors \(B_P,D_f,D_a\)
are pairwise coprime: the last two are \(1\) modulo \(B_P\), while
\[
 Q\equiv-\Delta,\quad \delta\equiv-A^2,\quad
 V\equiv-c,\quad C_y\equiv-D_m,\quad D_a\equiv A^2D_m\pmod{D_f}.
\]
Since \(D_f\equiv1\pmod A\) and \(\pmod{D_m}\), the claim follows.
The equation gives
\[
 (\delta Z/S)^2=B_PD_fD_a.
\]
Thus each absolute factor is a square. Since \(\Delta A^2\) is divisible
by four, \(D_f\equiv1\pmod4\), so it cannot be a negative square.
For \(c\ne0\), negativity of \(B_P\) would make \(D_f<0\); for
\(c=0\), the equation directly gives \(Z^2=S^2B_P\). Hence
\(B_P,D_f,D_a>0\). Write their positive roots as \(p_0,f,e\).
Then \(P>0\) forces \(B\ge2\), yielding the centered strip and authentic
packing bounds. The free guard additionally forces \(|Y|>|q|^3\).

A crucial order-of-proof point is that carrier integrality comes
\emph{before} main-root folding. If \(c=0\), it follows immediately from
\(Z^2=S^2B_P\). Otherwise set \(I=Up_0\) and \(T_o=p_0T\).
Then \(Q=(\Delta I)^2>1\), and the inherited adjacent-Pell denominator
descent makes \(y=e/\delta\) integral. Thus
\(Z^2=S^2p_0^2f^2y^2\), so in either branch \(S\mid A_c\).
This avoids assuming \(|A|\ge3\) while the weak scale might still be zero.

The valuation of the monic leading term in
\[
 A_c=z_c^5+3D_c^2z_c^3-19D_c^3z_c^2-57D_c^5
\]
then proves \(M\mid z_c\), and hence \(M\mid H(v-q^2)\) in the weak
case. Since \(H\equiv3\pmod q\), \(q\not\equiv0\pmod3\), and
\(H\) is odd, \(\gcd(H,M)=1\). Consequently
\[
 v=(2v_o+1)q^2,\qquad Y=(2v_o+1)q^3.
\]
Also \(H\mid K^5\). The restored magnitude and centered strip exclude
\(c=0\). Now the inherited early absolute-folding argument applies to
the strong and auxiliary norms; it yields \(D_m=d^2\) without using a
first Pell root. From \(p_0^2=D_mP\) it follows that \(P\) is a square,
and \(\gcd(D_t,D_g)=1\) splits it into \(D_t=\tau^2,D_g=g^2\).

All pre-projection signed-rank hypotheses are now present, including
\(c^2\mid I\), \(V\equiv-c\pmod f\), \(V\equiv-R\pmod c\), and
\(k=R+1+hXY\). The retained nonodd strong-rank and first-index arguments
give \(R=2t+1\), \(t\ge36\), and
\( |H|>|X|(|X|-3)^t\). Writing \(\sigma=\operatorname{sign}(c)\),
\[
 \Pi=(\sigma X-2^R)(2^R\sigma X+1),
\]
the signed projection gives \(H\mid\Pi^5\). The free-scale version first
gives only \(H\mid q^{15}\Pi^5\); its early carrier nevertheless gives
\(Y\) even. Both are covered by the stronger bound, with \(x=|X|\),
\[
 |q|^{15}|\Pi|^5<2^{20t-10}x^{25}<x^{t+1}/2^t<|H|,
\]
because \(x\ge2^{64}\) and
\(64(t-24)-(21t-10)=43t-1526\ge22\).
Thus \(\sigma X=2^R\); residues modulo three force \(\sigma=1\).
Only after this projection does the free case cancel its scale modulus,
since \(q\mid2^R\) and \(H\) is odd.

The inherited surviving signs and uniform centered half-binomial theorem
then give \(c,k,Y>0\) and \(0<\eta/k<1/4\). Fresh strong/auxiliary
private witnesses restore the literal old signed source zero, as detailed
in \texttt{PROOF.md}; no oddness of the incoming \(I/c^2\) is assumed.
Only this old zero is passed to the retained native/compiler/refutation
calibration.

\subsection{Completeness and its quantifier order}
A finite syntactic refutation supplies the canonical outer tuple with
\(D_m=d^2\), \(D_t=\tau^2\), \(A\) positive even,
\(R\equiv c\equiv3\pmod4\), and \(B>64\).
Since \(D_t\equiv1\pmod3\), \((3B-4)D_t\) is positive and
\(2\pmod3\), hence nonsquare. Fix a positive Pell guard pair \(s,g\),
and put \(p_0=d\tau g\). Choose one strong pair \(I,f\) with
\(4sAp_0c^2\mid I\), and set \(r=I/(Ap_0c^2)\). Keep it fixed.
Every prime of \(p_0sr\) divides \(I\), so the retained
supported-modulus interpolation theorem realizes every \(T_o\) residue
modulo \(p_0sr\) for this same pair.

The canonical projector gives \(\kappa=K/H\in\mathbb Z\), and
\[
 A_c/S=H^4\pi(T+\kappa+v_o)/(sr).
\]
Choose an intersective root \(a_0\) of \(\pi(T+\kappa+v_o)\) modulo
\(sr\), prescribe \(T_o\equiv p_0a_0\pmod{p_0sr}\), and use supported
interpolation. Then \(T=T_o/p_0\) and \(A_c/S\) are integers. Set
\(n=A_c/S+p_0fy-\lambda V\). The equation vanishes. The guard and
strong pair are never reselected after choosing the residue.

\subsection{Exact degree and replay boundary}
For the weak direct equation the register degrees are
\[
 \deg(B_P,U,D_f,Q,V,C_y,S,A_c)=(72,7,94,102,174,450,16,35).
\]
The unique dominating product \(S^2B_PD_fC_y\) has degree \(648\).
The attained monomial
\[
 j^{372}w^{64}v^{106}s^{16}r^{12}c^{34}z^{42}T^2
\]
has coefficient \(2^{546}3^{443}\beta^{372}>0\).
For the simple center,
\(\lambda^2-B_PD_f=U^2B_P^2-B_P\); exact expansion leaves a unique
leading term \(S^2QU^2B_P^2V^2\), of degree \(640\).
The free analogues have degrees \(580\) and \(574\).

The independent standard-library checker does not import the builder.
It binds every live arithmetic gate in free polynomial rings, checks
all identities and exact leading monomials, reconstructs every fixed
numeral, authenticates retained sources, and rejects every single-gate
mutation. Normal and optimized Python runs agree. These finite checks
are not substituted for unbounded Pell, interpolation, or compiler
theorems. The result is an ordinary ZFC-provable equivalence between
universal nonvanishing and syntactic \(\operatorname{Con}(\mathrm{ZFC})\),
not a proof of consistency or a new end-to-end formal certificate.
